\documentclass[10pt,a4paper]{amsart}
\usepackage{amsmath,amssymb,amsthm,graphicx}
\usepackage{verbatim}
\newenvironment{reference}{\comment}{\endcomment}
%\newenvironment{reference}{}{}
\newenvironment{slogan}{\comment}{\endcomment}
\newenvironment{history}{\comment}{\endcomment}

% For commutative diagrams we use Xy-pic
\usepackage[all]{xy}

% We use 2cell for 2-commutative diagrams.
\xyoption{2cell}
% Theorem environments.
%
\theoremstyle{plain}
\newtheorem{theorem}[subsection]{Theorem}
\newtheorem{proposition}[subsection]{Proposition}
\newtheorem{lemma}[subsection]{Lemma}

\theoremstyle{definition}
\newtheorem{definition}[subsection]{Definition}
\newtheorem{example}[subsection]{Example}
\newtheorem{exercise}[subsection]{Exercise}
\newtheorem{situation}[subsection]{Situation}

\theoremstyle{remark}
\newtheorem{remark}[subsection]{Remark}
\newtheorem{remarks}[subsection]{Remarks}

\numberwithin{equation}{subsection}

% Macros
%
\def\lim{\mathop{\mathrm{lim}}\nolimits}
\def\colim{\mathop{\mathrm{colim}}\nolimits}
\def\Spec{\mathop{\mathrm{Spec}}}
\def\Hom{\mathop{\mathrm{Hom}}\nolimits}
\def\Ext{\mathop{\mathrm{Ext}}\nolimits}
\def\SheafHom{\mathop{\mathcal{H}\!\mathit{om}}\nolimits}
\def\SheafExt{\mathop{\mathcal{E}\!\mathit{xt}}\nolimits}
\def\Sch{\mathit{Sch}}
\def\Mor{\mathop{\mathrm{Mor}}\nolimits}
\def\Ob{\mathop{\mathrm{Ob}}\nolimits}
\def\Sh{\mathop{\mathit{Sh}}\nolimits}
\def\NL{\mathop{N\!L}\nolimits}
\def\CH{\mathop{\mathrm{CH}}\nolimits}
\def\proetale{{pro\text{-}\acute{e}tale}}
\def\etale{{\acute{e}tale}}
\def\QCoh{\mathit{QCoh}}
\def\Ker{\mathop{\mathrm{Ker}}}
\def\Im{\mathop{\mathrm{Im}}}
\def\Coker{\mathop{\mathrm{Coker}}}
\def\Coim{\mathop{\mathrm{Coim}}}

% Boxtimes
%
\DeclareMathSymbol{\boxtimes}{\mathbin}{AMSa}{"02}

%
% Macros for moduli stacks/spaces
%
\def\QCohstack{\mathcal{QC}\!\mathit{oh}}
\def\Cohstack{\mathcal{C}\!\mathit{oh}}
\def\Spacesstack{\mathcal{S}\!\mathit{paces}}
\def\Quotfunctor{\mathrm{Quot}}
\def\Hilbfunctor{\mathrm{Hilb}}
\def\Curvesstack{\mathcal{C}\!\mathit{urves}}
\def\Polarizedstack{\mathcal{P}\!\mathit{olarized}}
\def\Complexesstack{\mathcal{C}\!\mathit{omplexes}}
% \Pic is the operator that assigns to X its picard group, usage \Pic(X)
% \Picardstack_{X/B} denotes the Picard stack of X over B
% \Picardfunctor_{X/B} denotes the Picard functor of X over B
\def\Pic{\mathop{\mathrm{Pic}}\nolimits}
\def\Picardstack{\mathcal{P}\!\mathit{ic}}
\def\Picardfunctor{\mathrm{Pic}}
\def\Deformationcategory{\mathcal{D}\!\mathit{ef}}
\usepackage{cleveref}
\crefname{originallemma}{original lemma}{original lemmas}
\Crefname{originallemma}{Original lemma}{Original lemmas}
\makeatletter
\newlabel{algebra-lemma-completion-flat}{{00MB}{source}{}{sourcepointer-1}{}}
\newlabel{algebra-lemma-completion-flat@cref}{{[originallemma][0][]00MB}{[1][0][]source}{}{}{}}
\makeatother
\crefname{originallemma}{original lemma}{original lemmas}
\Crefname{originallemma}{Original lemma}{Original lemmas}
\makeatletter
\newlabel{algebra-lemma-completion-tensor}{{00MA}{source}{}{sourcepointer-2}{}}
\newlabel{algebra-lemma-completion-tensor@cref}{{[originallemma][0][]00MA}{[1][0][]source}{}{}{}}
\makeatother
\crefname{originalsection}{original section}{original sections}
\Crefname{originalsection}{Original section}{Original sections}
\makeatletter
\newlabel{algebra-section-netherlander}{{00S0}{source}{}{sourcepointer-3}{}}
\newlabel{algebra-section-netherlander@cref}{{[originalsection][0][]00S0}{[1][0][]source}{}{}{}}
\makeatother
\crefname{originalequation}{original equation}{original equations}
\Crefname{originalequation}{Original equation}{Original equations}
\makeatletter
\newlabel{equation-C-prime}{{0AL4}{source}{}{sourcepointer-4}{}}
\newlabel{equation-C-prime@cref}{{[originalequation][0][]0AL4}{[1][0][]source}{}{}{}}
\makeatother
\crefname{originalreference}{original reference}{original references}
\Crefname{originalreference}{Original reference}{Original references}
\makeatletter
\newlabel{lemma-topologically-finite-type-Noetherian}{{0AJQ}{source}{}{sourcepointer-5}{}}
\newlabel{lemma-topologically-finite-type-Noetherian@cref}{{[originalreference][0][]0AJQ}{[1][0][]source}{}{}{}}
\makeatother
\title{The completed naive cotangent complex is the completion of a chosen presentation complex}
\author{The Stacks Project Authors}
\date{stacks-project:restricted.tex:local-snapshot-20261012; GNU FDL 1.2 or later}
\begin{document}
\maketitle
\noindent\textit{Editorial note.} All delivered mathematical text is the author's own, copied from the pinned Stacks Project chapter file restricted.tex; no statement or proof step was rewritten, reordered or removed.\par\medskip
\noindent\textit{Editorial note.} The counted proof is the author's complete argument: the proof identifies the conormal term under completion through the exactness and tensor agreement of completion on finite modules, reads off the equality of complexes, and derives the derived version from the homotopy equivalence with NL(alpha) and flatness\par\medskip
\noindent\textit{Editorial note.} The label lemma-NL-is-completion names the counted statement, and the proof environment immediately following it in the source is delivered verbatim as its proof.\par\medskip
\noindent\textit{Editorial note.} The counted statement is carried by the author's own lemma environment in the chapter source.\par\medskip
\noindent\textit{Allowed standard prerequisites (editorial).}\begin{itemize}
\item B\textasciicircum{}wedge is an object of the completed category (lemma-topologically-finite-type-Noetherian of the same chapter)
\item I-adic completion is exact on finite modules over a polynomial ring and agrees with tensoring with the completed polynomial ring (Algebra, Lemmas algebra-lemma-completion-tensor and algebra-lemma-completion-flat)
\item flatness of A[x]\textasciicircum{}wedge over A[x] and of B\textasciicircum{}wedge over B
\item the naive cotangent complex NL(alpha) of a presentation J/J\textasciicircum{}2 to the free module of differentials (Algebra, Section algebra-section-netherlander)
\end{itemize}\medskip
\ifdefined\hypertarget\hypertarget{sourcepointer-1}{}\fi
\noindent\textit{Source reference.} Original lemma 00MB: sources/stacks-chow/tags.txt, source lines 817--817. Prerequisite pointer invoked at proof line(s) 579; it names a standard fact of a sibling chapter that the counted argument uses, and every step of the counted proof built on it is retained verbatim in the delivered body.\par\medskip
\ifdefined\hypertarget\hypertarget{sourcepointer-2}{}\fi
\noindent\textit{Source reference.} Original lemma 00MA: sources/stacks-chow/tags.txt, source lines 816--816. Prerequisite pointer invoked at proof line(s) 578; it names a standard fact of a sibling chapter that the counted argument uses, and every step of the counted proof built on it is retained verbatim in the delivered body.\par\medskip
\ifdefined\hypertarget\hypertarget{sourcepointer-3}{}\fi
\noindent\textit{Source reference.} Original section 00S0: sources/stacks-chow/tags.txt, source lines 1022--1022. Prerequisite pointer invoked at proof line(s) 587; it names a standard fact of a sibling chapter that the counted argument uses, and every step of the counted proof built on it is retained verbatim in the delivered body.\par\medskip
\ifdefined\hypertarget\hypertarget{sourcepointer-4}{}\fi
\noindent\textit{Source reference.} Original equation 0AL4: sources/stacks-restricted/pinned-original-source.tex, source lines 326--326. Prerequisite pointer invoked at proof line(s) 572; it names a standard fact of the same chapter that the counted argument uses, and every step of the counted proof built on it is retained verbatim in the delivered body.\par\medskip
\ifdefined\hypertarget\hypertarget{sourcepointer-5}{}\fi
\noindent\textit{Source reference.} Original reference 0AJQ: sources/stacks-restricted/pinned-original-source.tex, source lines 395--395. Prerequisite pointer invoked at proof line(s) 573; it names a standard fact of the same chapter that the counted argument uses, and every step of the counted proof built on it is retained verbatim in the delivered body.\par\medskip
\begin{lemma}
\label{lemma-NL-is-completion}
Let $A$ be a Noetherian ring and let $I \subset A$ be an ideal.
Let $A \to B$ be a finite type ring map. Choose a presentation
$\alpha : A[x_1, \ldots, x_n] \to B$. Then
$\NL_{B^\wedge/A}^\wedge = \lim \NL(\alpha) \otimes_B B^\wedge$
as complexes and
$\NL_{B^\wedge/A}^\wedge = \NL_{B/A} \otimes_B^\mathbf{L} B^\wedge$
in $D(B^\wedge)$.
\end{lemma}

\begin{proof}
The statement makes sense as $B^\wedge$ is an object of
(\ref{equation-C-prime}) by
Lemma \ref{lemma-topologically-finite-type-Noetherian}.
Let $J = \Ker(\alpha)$. The functor of taking
$I$-adic completion is exact on finite modules over
$A[x_1, \ldots, x_n]$ and agrees with the functor
$M \mapsto M \otimes_{A[x_1, \ldots, x_n]} A[x_1, \ldots, x_n]^\wedge$, see
Algebra, Lemmas \ref{algebra-lemma-completion-tensor} and
\ref{algebra-lemma-completion-flat}. Moreover, the ring maps
$A[x_1, \ldots, x_n] \to A[x_1, \ldots, x_n]^\wedge$
and $B \to B^\wedge$ are flat. Hence
$B^\wedge = A[x_1, \ldots, x_n]^\wedge / J^\wedge$ and
$$
(J/J^2) \otimes_B B^\wedge = (J/J^2)^\wedge = J^\wedge/(J^\wedge)^2
$$
Since $\NL(\alpha) = (J/J^2 \to \bigoplus B\text{d}x_i)$,
see Algebra, Section \ref{algebra-section-netherlander},
we conclude the complex $\NL_{B^\wedge/A}^\wedge$ is equal
to $\NL(\alpha) \otimes_B B^\wedge$. The final statement
follows as $\NL_{B/A}$ is homotopy equivalent to $\NL(\alpha)$
and because the ring map $B \to B^\wedge$ is flat (so derived
base change along $B \to B^\wedge$ is just base change).
\end{proof}

\end{document}
